\documentclass[11pt]{article}
\usepackage[a4paper,margin=1in]{geometry}
\usepackage{amsmath,amssymb,amsthm,mathtools}
\usepackage{graphicx}
\usepackage[hidelinks]{hyperref}
\usepackage{microtype}

\newtheorem{theorem}{Theorem}
\newtheorem{lemma}[theorem]{Lemma}
\newcommand{\Hq}{\mathbb H}
\newcommand{\C}{\mathbb C}
\newcommand{\R}{\mathbb R}
\newcommand{\conv}{\operatorname{conv}}
\newcommand{\diag}{\operatorname{diag}}

\title{The $2\times2$ upper-bild trichotomy fails\\
in every dimension $n\geq3$}
\author{Candidate solution packet for arXiv:2008.02850}
\date{13 August 2026}

\begin{document}
\maketitle

\begin{center}
\fbox{\parbox{0.9\textwidth}{\textbf{Status.} Candidate full negative answer,
likely valid, awaiting specialist review.  The source paper already treats
two regimes at size three; the new ingredient is the missing singular
positive-semidefinite regime, including its exact upper bild, followed by a
uniform closure of all three regimes for every $n\geq3$.}}
\end{center}

\section{Source question and answer}

Carvalho, Diogo, and Mendes ask after Corollary~3.11 of
\cite{CDM} (source PDF, p.~9) whether any part of its $2\times2$ description
can be generalized to $n>2$:

\begin{center}
\includegraphics[width=0.98\textwidth]{figures/open_problem_crop.png}
\end{center}

More precisely, for $A=H+Si$ they ask whether the upper bild can still be
simplified according as $S$ is positive definite, positive semidefinite, or
indefinite.  The three formulas displayed in Corollary~3.11 are
\begin{align*}
 B^+(A)&=\conv\{W_{\C}^+(A),W_{\C}^+(A^*)\} &&(S\text{ indefinite}),\\
 B^+(A)&=\conv\{W_{\C}^+(A),v\}\quad(v\in\R)
       &&(S\text{ positive definite}),\\
 B^+(A)&=W_{\C}^+(A) &&(S\text{ positive semidefinite}).
\end{align*}
Because the second and third hypotheses overlap in standard terminology, we
read the last branch in the intended trichotomy as singular positive
semidefinite.  The answer is negative in every regime and every dimension
$n\geq3$.

\section{Definitions}

For $A\in M_n(\C)\subset M_n(\Hq)$, let
\[
 W_{\Hq}(A)=\{x^*Ax:x\in\Hq^n,\ x^*x=1\}.
\]
Writing $\C^+=\{a+bi:a,b\in\R,\ b\geq0\}$, its upper bild is
$B^+(A)=W_{\Hq}(A)\cap\C^+$.  We similarly write
$W_{\C}^+(A)=W_{\C}(A)\cap\C^+$.  Quaternionic units satisfy
$i^2=j^2=k^2=ijk=-1$; in particular
\begin{equation}\label{eq:flip}
 j^*ij=-i,\qquad j^*(1+i)j=1-i.
\end{equation}

\section{Proof intuition}

The source's two examples exploit two copies of the eigenvalue $1+i$.
Quaternionic conjugation by $j$ reverses its imaginary part, so averaging the
two copies produces the real point $1$.  This point lies outside the complex
hulls proposed in the indefinite and positive-definite branches.

To expose the missing singular-PSD branch, retain those two copies and add a
zero eigenvalue.  The same average produces $1$, already outside the complex
segment $[0,1+i]$.  More is true.  If the first two squared coordinate norms
sum to $a$, their imaginary contributions have total norm at most $a$.
Conversely, splitting those weights as $(a+b)/2$ and $(a-b)/2$ in quaternionic
directions $1$ and $j$ produces $a+bi$.  This identifies the entire upper bild
as a triangle.  Repeating the final diagonal entry then gives examples in
every higher dimension without changing the obstruction.

\section{The uniform counterexamples}

\begin{theorem}\label{thm:main}
For every integer $n\geq3$, set
\begin{align*}
 A_n^0&=\diag(1+i,1+i,0,\ldots,0),\\
 A_n^+&=\diag(1+i,1+i,i,\ldots,i),\\
 A_n^{\pm}&=\diag(1+i,1+i,-i,\ldots,-i).
\end{align*}
Then:
\begin{enumerate}
\item $A_n^0=H+Si$ has singular positive-semidefinite $S$, and
\[
 W_{\C}^+(A_n^0)=[0,1+i],\qquad
 B^+(A_n^0)=\conv\{0,1,1+i\}.
\]
In particular $B^+(A_n^0)\ne W_{\C}^+(A_n^0)$.
\item $A_n^+=H+Si$ has $S=I_n>0$, but there is no $v\in\R$ such that
\[
 B^+(A_n^+)=\conv\{W_{\C}^+(A_n^+),v\}.
\]
\item $A_n^{\pm}=H+Si$ has indefinite $S$, but
\[
 B^+(A_n^{\pm})\ne
 \conv\{W_{\C}^+(A_n^{\pm}),W_{\C}^+((A_n^{\pm})^*)\}.
\]
\end{enumerate}
Consequently none of the three formulas in Corollary~3.11 extends
universally to any dimension $n\geq3$.
\end{theorem}

The first assertion follows from an exact calculation.

\begin{lemma}\label{lem:triangle}
For $n\geq3$,
\[
 B^+(A_n^0)=\{a+bi:0\leq b\leq a\leq1\}
             =\conv\{0,1,1+i\}.
\]
\end{lemma}

\begin{proof}
Let $q=(q_1,\ldots,q_n)\in\Hq^n$ be a unit vector.  Write
$q_\ell=r_\ell u_\ell$ for $\ell=1,2$, where $r_\ell\geq0$ and
$u_\ell$ is a unit quaternion (with an arbitrary choice if $r_\ell=0$).
Only the first two diagonal entries contribute, so
\[
 q^*A_n^0q
 =a+r_1^2u_1^*iu_1+r_2^2u_2^*iu_2,
 \qquad a=r_1^2+r_2^2\in[0,1].
\]
Each $u_\ell^*iu_\ell$ is a unit pure imaginary quaternion.  Hence the
norm of the imaginary part is at most $r_1^2+r_2^2=a$.  If the displayed
value lies in $\C^+$, it has the form $a+bi$ with $0\leq b\leq a$.

Conversely, take $0\leq b\leq a\leq1$ and define
\[
 q=\left(\sqrt{\frac{a+b}{2}},\
          \sqrt{\frac{a-b}{2}}\,j,\
          \sqrt{1-a},0,\ldots,0\right).
\]
This is a unit vector, and \eqref{eq:flip} gives
\begin{align*}
q^*A_n^0q
&=\frac{a+b}{2}(1+i)+\frac{a-b}{2}(1-i)\\
&=a+bi.
\end{align*}
Thus every point in the triangle is attained.
\end{proof}

\begin{proof}[Proof of Theorem~\ref{thm:main}]
The matrix $A_n^0$ is normal, so its complex numerical range is the convex
hull of its eigenvalues, namely $[0,1+i]$.  Lemma~\ref{lem:triangle} proves
item 1, including strictness because $1$ belongs to the triangle but not to
the diagonal segment.

For $A_n^+$, the imaginary-part matrix is $I_n$, and normality gives
\[
 W_{\C}^+(A_n^+)=W_{\C}(A_n^+)=[i,1+i].
\]
Using only the first three coordinates and \eqref{eq:flip}, the unit vectors
\[
 q=2^{-1/2}(1,j,0,\ldots,0),\qquad
 p=2^{-1/2}(1,0,j,0,\ldots,0)
\]
give
\[
 q^*A_n^+q=1,\qquad p^*A_n^+p=\frac12.
\]
Thus $B^+(A_n^+)\cap\R$ contains two distinct points.  On the other hand,
for any real $v$, every point of $\conv\{[i,1+i],v\}$ is
$t z+(1-t)v$ with $z\in[i,1+i]$ and $0\leq t\leq1$; its imaginary part is
$t$.  Therefore that hull meets $\R$ only at $v$.  Item 2 follows.

Finally, the imaginary-part matrix of $A_n^{\pm}$ has both signs.  Its
complex numerical range and that of its adjoint are respectively
\[
 [-i,1+i]\quad\hbox{and}\quad[i,1-i].
\]
Taking upper halves yields
\[
 \conv\{W_{\C}^+(A_n^{\pm}),W_{\C}^+((A_n^{\pm})^*)\}
 =\conv\{i,\tfrac12,1+i\}.
\]
This triangle meets the real axis only at $1/2$.  Yet the same vector
$q=2^{-1/2}(1,j,0,\ldots,0)$ gives
$q^*A_n^{\pm}q=1$.  Hence $1$ lies in the actual upper bild and outside the
proposed hull, proving item 3.
\end{proof}

\section{Verification, scope, and novelty}

The exact quaternion identities and the symbolic parametrization in
Lemma~\ref{lem:triangle} are checked in
\texttt{code/verify\_counterexamples.py}.  That script also performs an
independent random-vector check of the norm bound.  The latter is a sanity
check only; the proof is the triangle inequality and explicit construction
above.

The claim answers the source question in its intended concrete sense: none of
the three formulas in the $2\times2$ trichotomy is a universal formula in any
higher dimension.  It does not rule out other, differently formulated
simplifications.  Example~3.12 of the source already supplies the $n=3$
indefinite and positive-definite failures.  The new mathematical content is
the singular-PSD failure, its exact upper bild, and the dimension-uniform
three-family closure.

A bounded novelty search through 13 August 2026 used the source title and
arXiv id, exact matrix and formula searches, general web search, the run
registry, and local title/abstract/full-source indexes.  It found no separate
occurrence of the singular-PSD triangle or the full uniform statement.  This
is not an exhaustive priority search.

\section{Human-review recommendation}

The result is elementary and likely valid.  Review should begin with
Lemma~\ref{lem:triangle}, especially the assertion that conjugation of $i$ by
a unit quaternion preserves the norm of the pure imaginary part and the
explicit converse vector.  Novelty should be judged as a focused completion:
two branches were already known from the source paper.

\begin{thebibliography}{9}
\bibitem{CDM}
L. Carvalho, C. Diogo, and S. Mendes,
\emph{Quaternionic numerical range of complex matrices},
Linear Algebra Appl. \textbf{620} (2021), 168--181;
arXiv:2008.02850.
\end{thebibliography}

\end{document}
